Throughout, $\operatorname{Cl}(\cdot)$ is the ideal class group (finite abelian), $\operatorname{Cl}(\cdot)_\ell$ its $\ell$-primary part (the subgroup of elements of $\ell$-power order), $\iota:\operatorname{Cl}(K)\to\operatorname{Cl}(L)$ the map induced by extension of ideals $\mathfrak a\mapsto\mathfrak a\mathcal O_L$, and $N:\operatorname{Cl}(L)\to\operatorname{Cl}(K)$ the map induced by the ideal norm $N_{L/K}$. Both are group homomorphisms; both restrict to the $\ell$-parts because a homomorphism of finite abelian groups maps elements of $\ell$-power order to elements of $\ell$-power order.

\emph{Step 1: $N\circ\iota$ is multiplication by $[L:K]$.} For an ideal $\mathfrak a$ of $K$, $N_{L/K}(\mathfrak a\mathcal O_L)=\mathfrak a^{[L:K]}$ (transitivity of the norm: the norm of an extended ideal is the $[L:K]$-th power). Hence $N\circ\iota=[L:K]=p^k$ on $\operatorname{Cl}(K)$.

\emph{Step 2: injectivity on the $\ell$-part.} Multiplication by $p^k$ is a bijection of the finite abelian $\ell$-group $\operatorname{Cl}(K)_\ell$ because $\gcd(p^k,|\operatorname{Cl}(K)_\ell|)=1$ ($\ell\ne p$); an inverse is multiplication by an integer $u$ with $up^k\equiv1$ modulo the exponent of $\operatorname{Cl}(K)_\ell$. Since $N\circ\iota$ is bijective on $\operatorname{Cl}(K)_\ell$, $\iota$ is injective there, and $\iota(\operatorname{Cl}(K)_\ell)\cap\ker N=0$: an element $\iota y$ with $y\in\operatorname{Cl}(K)_\ell$ and $N\iota y=0$ has $p^ky=0$, hence $y=0$.

\emph{Step 3: the decomposition.} Let $x\in\operatorname{Cl}(L)_\ell$. Then $N(x)\in\operatorname{Cl}(K)_\ell$, and we may put $y:=u\,N(x)$ with $u$ as in Step~2, so that $p^ky=N(x)$. Then $N(x-\iota y)=N(x)-N\iota y=N(x)-p^ky=0$, i.e.\ $x-\iota y\in\ker N$; it lies in $\operatorname{Cl}(L)_\ell$ as a difference of two elements of $\ell$-power order. Thus $x=\iota y+(x-\iota y)$ exhibits $\operatorname{Cl}(L)_\ell=\iota(\operatorname{Cl}(K)_\ell)+\ker(N)_\ell$, and the sum is direct by Step~2. Consequently
\[ |\operatorname{Cl}(L)_\ell|=|\operatorname{Cl}(K)_\ell|\cdot|\ker(N)_\ell| . \]

\emph{Step 4: the valuation statement.} For a finite abelian group $G$, $|G_\ell|=\ell^{v_\ell(|G|)}$, so $v_\ell(h_L)$ equals $v_\ell$ of $|\operatorname{Cl}(L)_\ell|$, and likewise for $K$. Taking $v_\ell$ of the displayed identity, $v_\ell(h_L)-v_\ell(h_K)=v_\ell(|\ker(N)_\ell|)\ge0$, with equality iff $\ker(N)_\ell=0$. Hence $v_\ell(h_L)>v_\ell(h_K)$ iff $\ker(N)_\ell\ne0$.

\emph{Step 5: the Weber tower.} Apply this with $K=B_{n-1}$, $L=B_n$, $p=2$, $k=1$: $B_n/B_{n-1}$ is cyclic of degree $2$ (Galois group $\langle\tau\rangle$). Then $v_\ell(h_n)\ge v_\ell(h_{n-1})$ for every odd $\ell$, which is the odd part of the integrality $h_{n-1}\mid h_n$ (\cite{KY} Eq.~(17): $k_n=h_n/h_{n-1}\in\Z$ and $k_n=[\RE:A_n]$), and
\[ \ell\mid k_n\iff v_\ell(h_n)>v_\ell(h_{n-1})\iff\ker(N)_\ell\ne0 : \]
the odd primes dividing $k_n$ are exactly the odd primes of the ``new'' part $\ker N$ of the class group. This is the sense in which $k_n$, and not $h_n$, is the natural object at a fixed layer; \cite{KY} Eq.~(17) is the description of $k_n$ by relative units that Theorem~A uses. Nothing about $\ell=p=2$ is claimed.
